\documentclass{article}
\usepackage[UTF8]{ctex}
\usepackage{amsmath,amssymb,amsfonts}
\usepackage{enumitem}
\usepackage{geometry}
\geometry{a4paper, margin=1in}
% === 错题本排版优化（不参与题目标识） ===
\setlength{\parindent}{0pt}
\setlength{\parskip}{0.35em plus 0.10em minus 0.08em}
\setlength{\emergencystretch}{3em}
\setlist{itemsep=0.25em, topsep=0.40em, parsep=0pt}
\allowdisplaybreaks[2]
% === 错题本排版优化结束 ===
\begin{document}
\section*{重庆大学2025年数学分析真题}
\section*{63．重庆大学2025年数学分析试卷}

\subsection{题目1}
一．（20 分）计算函数极限 \(\lim _{x \rightarrow+\infty} x\left[\frac{1}{e}-\left(\frac{x}{1+x}\right)^{x}\right]\) ．

\subsection{题目2}
二．（20分）设 \(S\) 为 \(\mathbb{R}^{3}\) 中的光滑封闭曲面，取外侧，考虑第二类曲面积分

\[
I=\iint_{S}\left(x^{3}+y^{2025}\right) \mathrm{d} y \mathrm{~d} z+\left(2 y^{3}+y+z^{2024}\right) \mathrm{d} z \mathrm{~d} x+\left(3 z^{3}-4 z-1929 x\right) \mathrm{d} x \mathrm{~d} y
\]

（1）试确定曲面 \(S\) 的方程，使得积分 \(I\) 的值最小，并求出这个最小值。\\
（2）将第（1）中得到的曲面 \(S\) 在第一卦限的部分记作 \(S_{1}\) ，求 \(S_{1}\) 的切平面，使得该切平面与三个坐标轴围成的几何体的体积最小。

\subsection{题目3}
三。（20 分）解答如下问题：\\
（1）已知数列 \(\left\{x_{n}\right\}_{n=1}^{\infty}\) 满足 \(\left|x_{n+1}-x_{n}\right|<r\left|x_{n}-x_{n-1}\right|, n=2,3, \cdots\) ．其中 \(0<r<1\) 为某个常数，证明：数列 \(\left\{x_{n}\right\}_{n=1}^{\infty}\) 收敛。\\
（2）如果将第（1）小题中的常数 \(r\) 换成 1 ，是否仍有 \(\left\{x_{n}\right\}_{n=1}^{\infty}\) 收敛？证明你的结论。

\subsection{题目4}
四．（10 分）设函数 \(f(x)\) 在 \([0,1]\) 上具有连续的一阶导函数．\\
（1）证明：对任意的 \(x \in[0,1]\) ，有

\[
|f(x)| \leq \int_{0}^{1}|f(t)| \mathrm{d} t+\int_{0}^{1}\left|f^{\prime}(t)\right| \mathrm{d} t
\]

（2）证明不等式：

\[
\left|f\left(\frac{1}{2}\right)\right| \leq \int_{0}^{1}|f(t)| \mathrm{d} t+\frac{1}{2} \int_{0}^{1}\left|f^{\prime}(t)\right| \mathrm{d} t
\]

\subsection{题目5}
五．（15 分）判断下列命题的正误，正确的需给出证明，错误的要举出反例．\\
（1）若级数 \(\sum_{n=1}^{\infty} x_{n}\) 收敛，则级数 \(\sum_{n=1}^{\infty}(-1)^{n} \frac{x_{n}}{n}\) 也收敛。\\
（2）若连续函数列 \(\left\{f_{n}(x)\right\}_{n=1}^{\infty}\) 在 \([0,1]\) 上逐点收敛于连续函数 \(f(x)\) ，则 \(\left\{f_{n}(x)\right\}_{n=1}^{\infty}\) 在 \([0,1]\) 上一致收敛于 \(f(x)\) 。\\
（3）若函数 \(f(x), g(x)\) 都在 \([0,+\infty)\) 上一致连续，且 \(g(x)\) 以 \(x\) 轴为水平渐进线，则函数 \(f(x) g(x)\)在 \([0,+\infty)\) 上一致连续．

\subsection{题目6}
六。（20 分）设 \(D\) 是 \(\mathbb{R}^{2}\) 中由光滑简单封闭曲线 \(C\) 所围成的闭区域，二元函数 \(f\) 和 \(g\) 在 \(D\) 上具有连续的二阶偏导数，记 \(\Delta f=\frac{\partial^{2} f}{\partial x^{2}}+\frac{\partial^{2} f}{\partial y^{2}}\) ．\\
（1）证明等式：

\[
\iint_{D}(g \Delta f) \mathrm{d} x \mathrm{~d} y=\int_{C} g \frac{\partial f}{\partial n} \mathrm{~d} s-\iint_{D}[\operatorname{grad}(f) \cdot \operatorname{grad}(g)] \mathrm{d} x \mathrm{~d} y .
\]

其中 \(n\) 为曲线 \(C\) 的外法单位向量，\(\frac{\partial f}{\partial n}\) 表示 \(f\) 沿 \(n\) 方向的方向导数，向量

\[
\operatorname{grad}(f)=\left(\frac{\partial f}{\partial x}, \frac{\partial f}{\partial g}\right), \operatorname{grad}(g)=\left(\frac{\partial g}{\partial x}, \frac{\partial g}{\partial g}\right) .
\]

（2）证明：若 \(f\) 在曲线 \(C\) 上满足 \(f \equiv 0\) ，则

\[
\iint_{D}\left(|f|^{2}+|\Delta f|^{2}\right) \mathrm{d} x \mathrm{~d} y \geq 2 \iint_{D}|\operatorname{grad}(f)|^{2} \mathrm{~d} x \mathrm{~d} y .
\]

\subsection{题目7}
七．（ 25 分）解答如下问题：

\begin{itemize}
  \item[（1）] 证明：含参变量 \(y\) 的广义积分 \(\int_{0}^{+\infty} e^{-x y} \frac{\sin x}{x} \mathrm{~d} x\) 在 \([0,+\infty)\) 上一致收敛。
  \item[（2）] 证明：在 \((0,+\infty)\) 上，含参变量 \(y\) 的广义积分 \(\int_{0}^{+\infty} e^{-x y} \sin x \mathrm{~d} x\) 不一致收敛，但内闭一致收敛。
  \item[（3）] 利用第（1）和第（2）小题中的结论计算 Dirichlet 积分 \(\int_{0}^{+\infty} \frac{\sin x}{x} \mathrm{~d} x\) 。
  \item[（4）] 利用第（3）小题中的结论证明： \(\int_{0}^{+\infty}\left[\int_{0}^{+\infty} \frac{\sin x \sin (x+y)}{x(x+y)} \mathrm{d} y\right] \mathrm{d} x=\frac{\pi^{2}}{8}\) ．
\end{itemize}

\subsection{题目8}
八．（20 分）设 \(f(x)\) 为 \([-\pi, \pi]\) 上的实值 Riemann 可积函数．给定区间 \(I \subset[-\pi, \pi], I\) 上的特征函数定义为

\[
\chi_{I}(x)= \begin{cases}1, & x \in I \\ 0, & x \notin I\end{cases}
\]

如果 \(I_{1}, I_{2}, \cdots, I_{N}\) 为 \([-\pi, \pi]\) 中两两不交的子空间，\(a_{1}, a_{2}, \cdots, a_{N}\) 均为常数，且 \(\bigcup_{k=1}^{N} I_{k}=[-\pi, \pi]\) ，那么我们称有限线性组合

\[
a_{1} \chi_{I_{1}}+a_{2} \chi_{I_{2}}+\cdots+a_{N} \chi_{I_{N}}
\]

为 \([-\pi, \pi]\) 上的阶梯函数．

\begin{itemize}
  \item[（1）] 证明：对任意的 \(\varepsilon>0\) ，存在 \([-\pi, \pi]\) 上的阶梯函数 \(\varphi(x)\) 和 \(\psi(x)\) 使得
\[
\varphi(x) \leq f(x) \leq \psi(x), x \in[-\pi, \pi] .
\]
并且
\[
\int_{-\pi}^{\pi}|\varphi(x)-f(x)| \mathrm{d} x<\varepsilon, \int_{-\pi}^{\pi}|\psi(x)-f(x)| \mathrm{d} x<\varepsilon .
\]
  \item[（2）] 利用（1）中的结论证明
\[
\lim _{p \rightarrow+\infty} \int_{-\pi}^{\pi} f(x) \sin (p x) \mathrm{d} x=\lim _{p \rightarrow+\infty} \int_{-\pi}^{\pi} f(x) \cos (p x) \mathrm{d} x=0
\]
  \item[（3）] 将 \(f(x)\) 延拓为 \((-\infty,+\infty)\) 上以 \(2 \pi\) 为周期的周期函数，延拓之后的函数仍记为 \(f(x)\) 。证明：\(f(x)\)的 Fourier 级数的部分和 \(S_{n}(x)\) 可写成
\[
S_{n}(x)=\frac{1}{\pi} \int^{\pi} \frac{f(x+t)}{2 \sin \underline{t}} \sin \left(\frac{2 n+1}{2} t\right) \mathrm{d} t, n=1,2, \cdots .
\]
\end{itemize}
\end{document}
